\documentclass[11pt,a4paper]{article}
\usepackage[margin=25mm,headheight=15pt]{geometry}
\usepackage[T1]{fontenc}
\usepackage[utf8]{inputenc}
\usepackage{lmodern,amsmath,amssymb,amsthm,mathtools,bm}
\usepackage{xcolor,graphicx,booktabs,tabularx,array,longtable}
\usepackage{microtype,enumitem,fancyhdr}
\usepackage{xurl}
\usepackage{hyperref}
\definecolor{accent}{HTML}{163B55}
\hypersetup{colorlinks=true,linkcolor=accent,citecolor=accent,urlcolor=accent,
 pdftitle={Generic backward uniqueness for linear retarded equations},pdfauthor={Matthew J. Colbrook},
 pdfsubject={AIM 921: mathematical result and proof}}
\urlstyle{same}
\setlength{\parindent}{0pt}
\setlength{\parskip}{0.35em}
\setlength{\emergencystretch}{3em}
\setlist{itemsep=0.25em,topsep=0.3em}
\newtheorem{theorem}{Theorem}[section]
\newtheorem{lemma}[theorem]{Lemma}
\newtheorem{proposition}[theorem]{Proposition}
\newtheorem{corollary}[theorem]{Corollary}
\theoremstyle{remark}
\newtheorem{remark}[theorem]{Remark}
\newtheorem{definition}[theorem]{Definition}
\DeclareMathOperator{\Tr}{Tr}
\DeclareMathOperator{\supp}{supp}
\DeclareMathOperator{\Lip}{Lip}
\DeclareMathOperator{\Rea}{Re}
\newcommand{\R}{\mathbb R}
\newcommand{\C}{\mathbb C}
\newcommand{\E}{\mathbb E}
\newcommand{\N}{\mathbb N}
\newcommand{\ip}[2]{\langle #1,#2\rangle}
\pagestyle{fancy}
\fancyhf{}
\fancyhead[L]{\small\sffamily AIM 921}
\fancyhead[R]{\small\sffamily Matthew J. Colbrook}
\fancyfoot[C]{\thepage}
\renewcommand{\headrulewidth}{0.3pt}
\allowdisplaybreaks[2]
\title{\textbf{Generic backward uniqueness for linear retarded equations}}
\author{Matthew J. Colbrook\\[3pt]
{\small Department of Applied Mathematics and Theoretical Physics}\\
{\small University of Cambridge, United Kingdom}\\
{\small\href{mailto:mjc249@maths.cam.ac.uk}{\texttt{mjc249@maths.cam.ac.uk}}}}
\date{3 November 2026}
\begin{document}
\maketitle
\begin{abstract}
We prove residual backward uniqueness for all pairs of real times in the specified compact C1 coefficient topology. A continuous evolution on a Hilbert phase space makes injectivity a countable intersection of open conditions, while a dense gap family supplies the one-to-one solution maps.
\end{abstract}
\noindent\textbf{Repository target.}
This manuscript concerns AIM problem 921 at revision
\texttt{5d7e26e8444e}. The precise scope is stated below.
\medskip

\section{Main result and proof}
\label{sec:921}

\paragraph{Answer.} The set in Problem~921 is residual. The useful change
of phase space is to retain the current value separately and put the past
history in $L^2$. The $C^1$ dependence of the history functional in operator
norm gives a common finite measure dominating all delay measures on a
compact time interval. This makes the enlarged Hilbert-space evolution
well defined and continuous in the coefficient. Injectivity is a $G_\delta$
condition for operators on a separable Hilbert space. The gap truncations
behind the density theorem of Hale--Oliva~\cite{HaleOliva1976} are injective
on this enlarged space as well.

Fix $r>0$ and an integer $n\ge1$, and use exactly the coefficient space
and topology in the question:
\[
 C_r=C([-r,0];\mathbb R^n),\qquad
 \mathcal B_r=\mathcal L(C_r,\mathbb R^n),\qquad
 \mathcal X=C^1(\mathbb R;\mathcal B_r).
\]
Introduce the separable real Hilbert space
\[
 \mathcal E=\mathbb R^n\oplus L^2([-r,0];\mathbb R^n).
 \tag{921.1}\label{eq:921-hilbert-phase}
\]
The two components of an initial state $(a,\phi)$ are independent: no
point value is assigned to the $L^2$ class $\phi$.

\begin{theorem}\label{thm:921}
There is a countable intersection $\mathcal R$ of open dense subsets of
$\mathcal X$ such that, for every $L\in\mathcal R$, its enlarged evolution
$U_L(t,\sigma)$ on $\mathcal E$ is injective for every real $\sigma<t$.
Consequently the set $\mathcal G_{r,n}$ in Problem~921 is residual in
$\mathcal X$.
\end{theorem}

\subsection{A common dominating measure}

By the Riesz representation theorem, write
\[
 L(t)f=\int_{[-r,0]}d\mu_t(\theta)f(\theta),
\]
where $\mu_t=(\mu_{t,ij})_{i,j=1}^n$ is a matrix of finite real signed
measures. The norm
$\|\mu\|_{\mathrm{var}}=\sum_{i,j}|\mu_{ij}|([-r,0])$ is equivalent,
with constants depending only on $n$, to the operator norm of $L$.
Thus $t\mapsto\mu_t$ is $C^1$ in variation norm.

Fix a compact time interval $I=[a,b]$. Define the finite positive measure
\[
 \nu=\sum_{i,j}\left(
 |\mu_{a,ij}|+\int_a^b|\mu'_{s,ij}|\,ds\right).
 \tag{921.2}\label{eq:921-envelope}
\]
The integral is an integral of finite measures in variation norm; the
variation map is continuous in that norm. The fundamental theorem of
calculus and the triangle inequality for variation give
\[
 \sum_{i,j}|\mu_{t,ij}|\le\nu\qquad(t\in I).
 \tag{921.3}\label{eq:921-domination}
\]
In particular, the matrix densities $R(t,\theta)=d\mu_t/d\nu$ may be
chosen jointly measurable, and
\[
 \|R(t,\theta)\|_{\mathrm{op}}
 \le\sum_{i,j}|R_{ij}(t,\theta)|\le1
 \quad\text{for almost every }(t,\theta)
 \text{ with respect to }dt\,d\nu(\theta).
\]
For justification of joint measurability, the densities form a continuous
map from $I$ to $L^1(\nu)$ entry by entry, because the corresponding
measures are continuous in variation norm. A jointly measurable
representative of a strongly measurable $L^1$-valued function suffices.

The same construction applies to a difference of two coefficient
families. On $[a,b]$, its envelope $\nu_\Delta$ satisfies
\[
 \nu_\Delta([-r,0])
 \le\|\mu_a-\widetilde\mu_a\|_{\mathrm{var}}
 +(b-a)\sup_{s\in[a,b]}
          \|\mu'_s-\widetilde\mu'_s\|_{\mathrm{var}}.
 \tag{921.4}\label{eq:921-difference-envelope}
\]
This is controlled by a constant times the compact $C^1$ distance in the
problem statement.

\subsection{Evolution on the Hilbert phase space}

Fix an initial time $\sigma$ and a terminal bound $b>\sigma$. Given
$(a,\phi)\in\mathcal E$, extend $\phi$ by zero outside $[-r,0]$ and define
the initial-history forcing, for almost every $s\in[\sigma,b]$, by
\[
 F_{L,\phi}(s)=\int_{[-r,0]}d\mu_s(\theta)\,
 \mathbf1_{\{s+\theta<\sigma\}}
 \phi(s+\theta-\sigma).
 \tag{921.5}\label{eq:921-forcing}
\]
Using an envelope from
\eqref{eq:921-envelope} on $[\sigma,b]$, Minkowski's inequality and
translation invariance of Lebesgue measure give
\[
 \begin{split}
 \|F_{L,\phi}\|_{L^2([\sigma,b])}
 &\le\int_{[-r,0]}
 \left(\int_\sigma^b
   \mathbf1_{\{s+\theta<\sigma\}}
   |\phi(s+\theta-\sigma)|^2\,ds\right)^{1/2}d\nu(\theta)\\
 &\le\nu([-r,0])\,\|\phi\|_2.
 \end{split}
 \tag{921.6}\label{eq:921-young}
\]
This also proves independence of the representative of $\phi$. Indeed,
translating a Lebesgue null set produces a null set for each fixed
$\theta$, and Fubini's theorem applies to $ds\,d\nu(\theta)$.

For a continuous future path $x:[\sigma,b]\to\mathbb R^n$, put
\[
 (B_Lx)(s)=\int_{[-r,0]}d\mu_s(\theta)\,
  \mathbf1_{\{s+\theta\ge\sigma\}}x(s+\theta).
\]
This is measurable and depends only on values of $x$ up to time $s$.
If $M=\sup_{s\in[\sigma,b]}\|\mu_s\|_{\mathrm{var}}$, then
\[
 |(B_Lx)(s)|\le M\sup_{\sigma\le u\le s}|x(u)|.
 \tag{921.7}\label{eq:921-volterra-bound}
\]
The Volterra integral equation
\[
 x(t)=a+\int_\sigma^tF_{L,\phi}(s)\,ds
          +\int_\sigma^t(B_Lx)(s)\,ds
 \tag{921.8}\label{eq:921-integral-evolution}
\]
therefore has a unique continuous solution, by Picard iteration on short
intervals or the usual factorial estimate for Volterra iteration.
It belongs to $W^{1,2}([\sigma,b];\mathbb R^n)$ and obeys, with
$T=b-\sigma$,
\[
 \sup_{\sigma\le t\le b}|x(t)|
 \le\bigl(|a|+\sqrt T\,\nu([-r,0])\|\phi\|_2\bigr)e^{MT}.
 \tag{921.9}\label{eq:921-path-bound}
\]

Concatenate the initial $L^2$ history and this future path, and define
\[
 U_L(t,\sigma)(a,\phi)=(x(t),x_t),
 \qquad x_t(\theta)=x(t+\theta)\quad\text{as an }L^2\text{ class}.
 \tag{921.10}\label{eq:921-evolution}
\]
Equations \eqref{eq:921-young}--\eqref{eq:921-path-bound} make this a
bounded linear operator on $\mathcal E$. A change of the concatenated
path at a single joining point has no effect on the equation almost
everywhere: this follows from the same Fubini argument used for
representative independence. Uniqueness then gives the evolution law
\[
 U_L(t,s)U_L(s,\sigma)=U_L(t,\sigma)
 \qquad(\sigma\le s\le t).
 \tag{921.11}\label{eq:921-evolution-law}
\]
For each initial state, $U_L(t,\sigma)(a,\phi)$ is strongly continuous
in $t\ge\sigma$, including at $t=\sigma$. At the initial time this uses
$x(t)\to a$, strong continuity of translations in $L^2$, and the fact
that the new future segment has length $t-\sigma\to0$ and bounded values.

\begin{lemma}\label{lem:921-parameter-continuity}
For each fixed pair $\sigma<t$, the map
$L\mapsto U_L(t,\sigma)$ from $\mathcal X$ into
$\mathcal L(\mathcal E)$ is continuous in operator norm.
\end{lemma}
\begin{proof}
Apply \eqref{eq:921-young} to the difference of two coefficient families:
\[
 \|F_{L,\phi}-F_{\widetilde L,\phi}\|_2
 \le\nu_\Delta([-r,0])\|\phi\|_2.
\]
For solutions with the same initial state, their difference $z=x-\widetilde x$
satisfies the difference of \eqref{eq:921-integral-evolution}, with
\[
 B_Lx-B_{\widetilde L}\widetilde x
 =B_Lz+(B_L-B_{\widetilde L})\widetilde x.
\]
The last term has absolute value at most
$\nu_\Delta([-r,0])\sup|\widetilde x|$.
Use \eqref{eq:921-path-bound} for $\widetilde x$ and Gronwall's inequality
for $z$. On a neighborhood of $L$ with a uniform bound in
$C^1([\sigma,t];\mathcal B_r)$, this gives
\[
 \sup_{\sigma\le s\le t}|z(s)|
 \le C\,\nu_\Delta([-r,0])\bigl(|a|+\|\phi\|_2\bigr),
\]
where $C$ is independent of the initial state and the nearby coefficient.
The difference of the terminal histories is at most $\sqrt r\sup|z|$
in $L^2$, since their older initial-history parts agree. Now use
\eqref{eq:921-difference-envelope} and take the supremum over the unit
ball of $\mathcal E$.
\end{proof}

\subsection{A dense family injective on the enlarged space}

Consider coefficients of the form
\[
 L(t)f=A(t)f(-r)+\int_{[-r+\rho,0]}d\eta_t(\theta)f(\theta),
 \qquad 0<\rho<r,
 \tag{921.12}\label{eq:921-gap-family}
\]
where $A(t)$ is invertible for almost every real $t$ and the coefficient
family is $C^1$ in variation norm.

\begin{lemma}\label{lem:921-gap-injective}
Every family \eqref{eq:921-gap-family} has injective $U_L(t,\sigma)$ on
$\mathcal E$ for every $\sigma<t$.
\end{lemma}
\begin{proof}
Suppose $U_L(t,\sigma)(a,\phi)=0$. Forward uniqueness extends the future
path by zero after $t$, so the concatenated path is zero almost everywhere
on $[t-r,\infty)$.

We show how to extend a zero region backward. Suppose that the path is
zero almost everywhere on $[b,\infty)$, with $b\ge\sigma-r$. For almost
every
\[
 s\ge\max\{\sigma,b+r-\rho\},
\]
its derivative is zero, and the integral in
\eqref{eq:921-gap-family} is zero because all arguments
$s+\theta\ge s-r+\rho\ge b$. The assertion about the integral is an
almost-everywhere assertion justified by a common dominating measure
and Fubini, exactly as in \eqref{eq:921-young}.
The equation gives $A(s)x(s-r)=0$. Since $A(s)$ is invertible almost
everywhere, the path is zero almost everywhere on
\[
 [\max\{\sigma-r,b-\rho\},\infty).
\]
Starting with $b=t-r$, finitely many iterations give zero on
$[\sigma-r,\infty)$. Thus $\phi=0$ in $L^2$. The future path is
continuous at $\sigma$, so its almost-everywhere vanishing also gives
$a=x(\sigma)=0$.
\end{proof}

\begin{lemma}\label{lem:921-density}
The families \eqref{eq:921-gap-family} are dense in $\mathcal X$.
\end{lemma}
\begin{proof}
Start with any coefficient $L$, represented by $\mu_t$, and any compact
time interval on which it is to be approximated. Write
$A(t)=\mu_t(\{-r\})$. For $\rho\downarrow0$, restriction to
$(-r,-r+\rho)$ tends strongly to zero on the space of finite measures
in variation norm. The restriction operators have norm at most one.
Their convergence is therefore uniform on each norm-compact set of
measures. Apply this to the compact sets of values of $\mu_t$ and
$\mu'_t$ on the chosen time interval. It follows that
\[
 \mu_t^{\rho}=A(t)\delta_{-r}+\mu_t|_{[-r+\rho,0]}
\]
approximates $\mu_t$ in compact $C^1$ norm as $\rho\downarrow0$.

It remains to perturb $A$ so that it is invertible almost everywhere.
For each fixed $t$, the polynomial
$S\mapsto\det(A(t)+S)$ on $M_n(\mathbb R)$ is not identically zero.
Its zero set has Lebesgue measure zero. Fubini's theorem on
$[-m,m]\times\{S:\|S\|<\epsilon\}$, followed by a countable intersection
over positive integers $m$, shows that for almost every arbitrarily
small constant matrix $S$,
\[
 \det(A(t)+S)\ne0\quad\text{for almost every }t\in\mathbb R.
\]
Replacing $A(t)$ by $A(t)+S$ changes the coefficient by
$f\mapsto Sf(-r)$ and does not change its time derivative. Combining
this perturbation with the gap truncation proves density in precisely
the required compact $C^1$ topology.
\end{proof}

\subsection{The Baire argument and the original continuous histories}

On a separable Hilbert space, the injective bounded operators form a
$G_\delta$ in operator norm. Here are the details needed in this setting.
Choose a countable dense set $D\subset\mathcal E$ and an orthonormal
basis $(e_j)$. For positive integers $j,m$, let
\[
 \mathcal O_{j,m}
 =\bigcup_{v\in D}
   \{T\in\mathcal L(\mathcal E):\|T^*v-e_j\|<1/m\}.
 \tag{921.13}\label{eq:921-open-tests}
\]
Each set is open. Moreover,
\[
 \bigcap_{j,m}\mathcal O_{j,m}
 =\{T:\overline{\operatorname{Ran}T^*}=\mathcal E\}
 =\{T:\ker T=\{0\}\}.
 \tag{921.14}\label{eq:921-injective-gdelta}
\]
The first equality uses density of $D$ and continuity of $T^*$; the
second is the Hilbert-space orthogonal-complement identity.

For every rational pair $\sigma<t$, take the preimages of the sets
\eqref{eq:921-open-tests} under $L\mapsto U_L(t,\sigma)$.
They are open by Lemma~\ref{lem:921-parameter-continuity}. Each contains
the dense gap family of Lemmas~\ref{lem:921-gap-injective} and
\ref{lem:921-density}, so each is dense. Their intersection over all
$j,m$ and all rational pairs $\sigma<t$ is a countable intersection
$\mathcal R$ of open dense subsets of $\mathcal X$. The latter is a
Fr\'echet space, hence a Baire space.

For $L\in\mathcal R$, injectivity holds for every real pair as well.
Suppose $U_L(t,\sigma)v=0$ with real $\sigma<t$. Choose a rational
$p>t$. For every rational $q\in(\sigma,t)$, the evolution law gives
\[
 U_L(p,q)U_L(q,\sigma)v
 =U_L(p,t)U_L(t,\sigma)v=0.
\]
The first factor is injective because $p,q$ are rational. Thus
$U_L(q,\sigma)v=0$. Let rational $q\downarrow\sigma$ and use strong
right continuity to conclude that $v=0$.

Finally, the bounded injective map
\[
 J:C_r\longrightarrow\mathcal E,\qquad
 J\phi=(\phi(0),\phi)
\]
identifies a continuous history with its compatible enlarged state.
The classical continuous-history solution satisfies
\eqref{eq:921-integral-evolution}, so uniqueness gives
\[
 U_L(t,\sigma)J\phi=J\bigl(T_L(t,\sigma)\phi\bigr).
\]
If $T_L(t,\sigma)\phi=0$, enlarged-space injectivity implies $J\phi=0$,
and hence $\phi=0$. Therefore $\mathcal R\subseteq\mathcal G_{r,n}$,
proving Theorem~\ref{thm:921}.

The proof uses the original $C^1$ coefficient topology throughout. It
does not infer openness of continuous-history injectivity, and it does
not replace distributed delays by finite step measures. The gap family
is used to prove density; the $G_\delta$ argument takes place on the
larger Hilbert phase space before returning to the specified history space.


\begin{thebibliography}{99}
\small
\setlength{\itemsep}{3pt}

\bibitem{HaleOliva1976}
J.~K.~Hale and W.~M.~Oliva,
\emph{One-to-oneness for linear retarded functional differential equations},
J. Differential Equations \textbf{20} (1976), 28--36.
\href{https://doi.org/10.1016/0022-0396(76)90093-0}{doi:10.1016/0022-0396(76)90093-0}.
\end{thebibliography}
\end{document}
